% Page budget: 1.50 pages | Owns: datasets, splits, excitation, preprocessing, comparators, metrics, and experimental protocol.
% WRITING GUIDE
% Purpose: Make the simulation and real-data evaluation reproducible and define fair comparisons before showing outcomes.
% Sources: README "Source map per subsection", rows the V rows (data-generating system, data, controller, compared models, metrics, hyperparameters); mandatory papers and the reference gate as in the README.
% Research-plan anchor: Preserve the planned interpolation, generalisation, and black-box questions; update datasets, horizons, excitation, and preprocessing to what was actually used.
% Current decisions: Independent phases, held-out positions, multisine-free motion, truth-only friction or absorber variants, matched budgets, and long-horizon evaluation.
% Choices to justify: Operating-point grid, ASMPT motion profiles, frequency bands, amplitudes, phase realizations, crest-factor selection, record length, sample rate, split sizes, window sampling, settling interval, primary metric, seed policy, comparator capacity, noise status, and all tuned hyperparameters.
% Hyperparameter check: Classify each important setting as physics-derived, inherited, validation-tuned, safety-constrained, or compute-constrained; give a sensitivity result for settings that could change the conclusion.
% Cautions: Plans are not evidence; use the corrected matched friction dataset only; do not call a generated profile ASMPT-relevant until its motion parameters are confirmed; keep oracle states outside the trainable pipeline.
% Planned figures: Reuse the training-horizon strip and show a 2D train, validation, and test operating-point coverage map; use 3D only if the third coordinate communicates additional coverage rather than visual clutter.
% Section is finished when: Every dataset, split, unit, horizon, comparator, metric, seed policy, and checkpoint rule is unambiguous.
% Header notes (2026-10-15): real-data evaluation and the settling interval are out of the thesis results
%   (THESIS-RESULTS.md Sec. 6); E5 (changed controller) is left out for now (Dirk 2026-10-15).
%   Dataset of every reported run: Thesis-writeup/Data/Coulomb-tanh-and-MSD-lowpass-noise, floor from
%   Coulomb-tanh-and-MSD-noise-free (runs.tsv rows 201 to 207, THESIS_NOISE=lowpass; D-238).
%   Data-based numbers are macros \Setup* from tables/setup_numbers.tex, written by the scripts in
%   Thesis-writeup/Code/Setup/ (see its README) on that dataset; numbers are never typed in.
% Header notes (2026-10-16, clean rewrite, handoff 2026-10-16-section-5-clean-rewrite):
%   Main text states only what Section VI uses; every design detail sits in Appendix app:setup with a
%   pointer from here (Jasper: pragmatic, details in an appendix). Five models in the main text (THESIS-RESULTS
%   "Presentation agreed 2026-10-16"). Justifications checked item by item (scratchpad reports
%   justification-check-data/-train, 2026-10-16); design choices may rest on the known simulated system
%   (D-259 item 1). Notation renames (Dirk 2026-10-16): coupling vector b_a, absorber displacement q_a,
%   friction scale v_fr, multisine force u^ms, multisine phases alpha, noise level lambda_j, black box f_bb,
%   h_bb, theta_bb, x^bb.
% MUST ESTABLISH (outline agreed with Dirk 2026-10-15, updated 2026-10-16)
%   1. The data come from a known system: the baseline plus a payload absorber and rail friction, both
%      absent from the baseline, in closed loop with one known controller K1 and with measurement noise at a
%      representative level (Sections II, III; Aspect 2).
%   2. The data excite what must be learned: motion profiles at low frequencies, a multisine over the
%      absorber band placed with the simulated frequency response (simulation only) (Aspect 2).
%   3. The test records separate interpolation, operational motion and range extrapolation; each E record
%      leaves the training range in one design setting against its I record; E6 is a stress record (Aspect 4).
%   4. Five compared models: M_LPV and M_LTI untrained, M_U and the final model M_OBC from a 10 % detuned
%      start, and the black box (no physics, sized on validation); choice-justifying runs in the appendix.
%   5. Metrics: NRMS_e per axis with the RMS in m and the record's noise floor; combination error and the
%      overlaps of Section IV with the predicted first-order bias (Aspect 3).
% GENERATED by Thesis-writeup/Code/Setup/*.py (see its README); never edit by hand.
% Section V (Experiment design) numbers computed on the thesis dataset.
% --- band_edge.py
\providecommand{\SetupCrossMax}{}\renewcommand{\SetupCrossMax}{106}% max K1 crossover [Hz]
\providecommand{\SetupCrossMin}{}\renewcommand{\SetupCrossMin}{75}% min K1 crossover over axes and training Y [Hz], simulated system
\providecommand{\SetupMotionFNN}{}\renewcommand{\SetupMotionFNN}{24}% max over motion-only training records and axes of the 99 % power frequency of u [Hz]
% --- band_frf.py
\providecommand{\SetupAddedMassDbHi}{}\renewcommand{\SetupAddedMassDbHi}{15}% highest closed-loop |S G|_YY of the system against the baseline, 100-300 Hz, training Y [dB]
\providecommand{\SetupAddedMassDbLo}{}\renewcommand{\SetupAddedMassDbLo}{-21}% lowest closed-loop |S G|_YY of the system against the baseline, 100-300 Hz, training Y [dB]
\providecommand{\SetupClPeak}{}\renewcommand{\SetupClPeak}{267}% closed-loop peak of T_YY with K1, median over training Y [Hz]
% --- controller_rate.py
\providecommand{\SetupKcheck}{}\renewcommand{\SetupKcheck}{1.5\cdot10^{-9}}% largest relative RMS, K1 at 20 kHz vs recorded u_fb
\providecommand{\SetupKholdMed}{}\renewcommand{\SetupKholdMed}{10.4}% median relative RMS [%] of the force change within one 4 kHz hold (block mean vs u_fb at kD)
\providecommand{\SetupKmismatchMax}{}\renewcommand{\SetupKmismatchMax}{1.0}% largest relative RMS [%], same
\providecommand{\SetupKmismatchMed}{}\renewcommand{\SetupKmismatchMed}{0.5}% median relative RMS [%], K1 at 4 kHz vs u_fb at the same instant
\providecommand{\SetupKmismatchMotion}{}\renewcommand{\SetupKmismatchMotion}{0.03}% largest relative RMS [%] in the records without multisine
% --- coverage.py
\providecommand{\SetupEAcc}{}\renewcommand{\SetupEAcc}{100}% E3 acceleration [% of max]
\providecommand{\SetupESixAcc}{}\renewcommand{\SetupESixAcc}{131}% E6 acceleration [% of max]
\providecommand{\SetupEVel}{}\renewcommand{\SetupEVel}{98}% E4 velocity [% of max]
\providecommand{\SetupEY}{}\renewcommand{\SetupEY}{0.39}% largest |Y| of E1, E2 [m]
\providecommand{\SetupTrainAcc}{}\renewcommand{\SetupTrainAcc}{75}% largest training acceleration [% of max]
\providecommand{\SetupTrainVel}{}\renewcommand{\SetupTrainVel}{75}% largest training velocity [% of max]
\providecommand{\SetupTrainY}{}\renewcommand{\SetupTrainY}{0.30}% largest |Y| of a training reference [m]
% --- filter_edge.py
\providecommand{\SetupFiltDecay}{}\renewcommand{\SetupFiltDecay}{2.6}% 1 % decay time of the anti-alias filter impulse response [ms]
\providecommand{\SetupFiltTau}{}\renewcommand{\SetupFiltTau}{0.5}% slowest pole time constant of the anti-alias filter [ms]
\providecommand{\SetupTrim}{}\renewcommand{\SetupTrim}{5}% samples removed at each record end [ms]
% --- filter_leak.py
\providecommand{\SetupAliasFilt}{}\renewcommand{\SetupAliasFilt}{4.0\cdot10^{-11}}% largest RMS of the noise above 2 kHz after the filter [m]
\providecommand{\SetupAliasRaw}{}\renewcommand{\SetupAliasRaw}{4.6\cdot10^{-9}}% largest RMS of the noise above 2 kHz, unfiltered [m]
\providecommand{\SetupFiltDist}{}\renewcommand{\SetupFiltDist}{2.3\cdot10^{-9}}% largest RMS change of the noise-free output by the rule [m]
% --- gap_snr.py
\providecommand{\SetupGapSnrMedMotion}{}\renewcommand{\SetupGapSnrMedMotion}{65}% median SNR 24-105 Hz, training records without multisine [dB]
\providecommand{\SetupGapSnrMedMs}{}\renewcommand{\SetupGapSnrMedMs}{56}% median SNR 24-105 Hz, training records with multisine [dB]
\providecommand{\SetupGapSnrMinMotion}{}\renewcommand{\SetupGapSnrMinMotion}{59}% min SNR 24-105 Hz, training records without multisine [dB]
\providecommand{\SetupGapSnrMinMs}{}\renewcommand{\SetupGapSnrMinMs}{49}% min SNR 24-105 Hz, training records with multisine [dB]
% --- loop_memory.py
\providecommand{\SetupDecay}{}\renewcommand{\SetupDecay}{44}% largest 1 % decay time of the PS impulse response [ms]
\providecommand{\SetupTauMax}{}\renewcommand{\SetupTauMax}{11}% slowest closed-loop time constant over training Y [ms]
\providecommand{\SetupWindow}{}\renewcommand{\SetupWindow}{100}% training window [ms]
% --- rail_velocity.py
\providecommand{\SetupFricVzero}{}\renewcommand{\SetupFricVzero}{1}% friction velocity scale 1/g [mm/s]
\providecommand{\SetupRailShareMed}{}\renewcommand{\SetupRailShareMed}{87}% median share of samples with |v| >= 3/g, multisine training records [%]
\providecommand{\SetupRailShareMin}{}\renewcommand{\SetupRailShareMin}{71}% smallest share of samples with |v| >= 3/g, multisine training records, any rail [%]
\providecommand{\SetupRailVrmsMin}{}\renewcommand{\SetupRailVrmsMin}{7.9}% smallest rail velocity RMS, multisine training records [mm/s]
% --- rate.py
\providecommand{\SetupPmGate}{}\renewcommand{\SetupPmGate}{5}% phase-margin gate [deg], D-166 HEURISTIC
\providecommand{\SetupPmLossFour}{}\renewcommand{\SetupPmLossFour}{3.8}% largest phase-margin loss at 4 kHz vs 20 kHz [deg], 5 training Y, 3 axes
\providecommand{\SetupPmLossTwo}{}\renewcommand{\SetupPmLossTwo}{8.6}% largest phase-margin loss at 2 kHz vs 20 kHz [deg]
% --- snr.py
\providecommand{\SetupNoiseRmsXone}{}\renewcommand{\SetupNoiseRmsXone}{1.1\cdot10^{-8}}% median RMS(y - y0) [m], axis X1
\providecommand{\SetupNoiseRmsXtwo}{}\renewcommand{\SetupNoiseRmsXtwo}{1.1\cdot10^{-8}}% median RMS(y - y0) [m], axis X2
\providecommand{\SetupNoiseRmsY}{}\renewcommand{\SetupNoiseRmsY}{5.8\cdot10^{-9}}% median RMS(y - y0) [m], axis Y
\providecommand{\SetupSnrMaxX}{}\renewcommand{\SetupSnrMaxX}{74}% max SNR [dB], both X rails
\providecommand{\SetupSnrMaxXone}{}\renewcommand{\SetupSnrMaxXone}{74}% max SNR [dB], axis X1
\providecommand{\SetupSnrMaxXtwo}{}\renewcommand{\SetupSnrMaxXtwo}{74}% max SNR [dB], axis X2
\providecommand{\SetupSnrMaxY}{}\renewcommand{\SetupSnrMaxY}{80}% max SNR [dB], axis Y
\providecommand{\SetupSnrMedXone}{}\renewcommand{\SetupSnrMedXone}{63}% median SNR [dB], axis X1
\providecommand{\SetupSnrMedXtwo}{}\renewcommand{\SetupSnrMedXtwo}{62}% median SNR [dB], axis X2
\providecommand{\SetupSnrMedY}{}\renewcommand{\SetupSnrMedY}{74}% median SNR [dB], axis Y
\providecommand{\SetupSnrMinX}{}\renewcommand{\SetupSnrMinX}{58}% min SNR [dB] over training records, both X rails
\providecommand{\SetupSnrMinXone}{}\renewcommand{\SetupSnrMinXone}{59}% min SNR [dB] over training records, axis X1
\providecommand{\SetupSnrMinXtwo}{}\renewcommand{\SetupSnrMinXtwo}{58}% min SNR [dB] over training records, axis X2
\providecommand{\SetupSnrMinY}{}\renewcommand{\SetupSnrMinY}{71}% min SNR [dB] over training records, axis Y
% generated by Thesis-writeup/Code/Setup (numbers of this section)
\section{Experiment design}\label{sec:setup}
\ifdraft
\begin{itemize}
\item Opening: system, records, compared models, metrics; design details and checks in Appendix~\ref{app:setup}
\end{itemize}
\fi
This section gives the data-generating system
(Section~\ref{sec:setup_system}), the records and their split
(Section~\ref{sec:setup_data}), the compared models
(Section~\ref{sec:setup_models}) and the metrics
(Section~\ref{sec:setup_metrics}).
The design details and their checks are in Appendix~\ref{app:setup}.

\subsection{Data-generating system}\label{sec:setup_system}
% Sources: D-204, D-224 (truth), D-188 (absorber), D-223, D-226 (friction gain), D-221 (K1),
% D-238 (noise); tdg_config.m, gtd_build_plant.m; Code/Setup rail_velocity.py, band_edge.py, snr.py.
\ifdraft
\begin{itemize}
\item The data are simulated: the system is the baseline plus rail friction and a payload absorber, both absent from the baseline
\item The friction subtracts from the actuator forces; a smoothed Coulomb law, Coulomb-like on the excited records
\item The absorber adds one coordinate, coupled through the inertia; matrices written out as in Section~II
\item Parameter values with their basis; both absorber features lie above the crossover of the loop
\item One known controller $K_1$, the same in every record; the injected force is the only feedforward
\item Noise matched to the machine; noise-free twin of every record; SNR (?: confidentiality)
\end{itemize}
\fi
The data are simulated, so the dynamics the baseline misses are known.
The data-generating system is the baseline \eqref{eq:eom} with Coulomb
friction on the three rails and a vibration absorber on the payload
(Fig.~\ref{fig:gantry_absorber}).
Both are absent from the baseline.
\begin{figure}[t]
  \centering
  \resizebox{\columnwidth}{!}{% ===========================================================================
% Final thesis figure: lumped-parameter model of the dual-drive H-type gantry
% with the simulated payload absorber. Geometry adapted from
% Garcia-Herreros et al. (2013), Fig. 2.
%
% Requires in the preamble:
%   \usepackage{tikz}
%   \usetikzlibrary{arrows.meta,calc,decorations.pathmorphing}
%
% Conventions
%   * Coordinates are in drawing units; the x/y keys of the tikzpicture set the
%     unit, so the drawing scales while the labels keep the document's own font
%     and size.  At 0.3mm the drawing is 1:1 with Garcia Fig. 2 (about 285pt
%     wide); 0.26mm fits one IEEE column (252pt).
%   * The origin is the cross arm centre of mass (the blue coordinate X).
%   * Parts on the cross arm are drawn in its rotated frame: u along the beam
%     (positive to the right), v perpendicular to it.  The measured coordinate
%     Y points LEFT in this view, so Y = -u, and the absorber, which sits at
%     +L0 in Y, is drawn to the LEFT of the payload.
%   * Changing \th (yaw) or \Yp (payload position) moves every attached part,
%     including its labels.  With \Yp > 0 the absorber ends up between the CoM
%     and the payload, and the F_Y label then needs a manual nudge.
%   * Schematic, not to scale: the absorber offset L0 and the mass sizes are
%     drawn for legibility.  The absorber rides on the payload, so it may
%     overhang the beam end.
%   * Colours: red = measured coordinates, blue = generalized coordinates,
%     grey = absorber (m_a, k_a, c_a, q_a) and rail friction F_fr, present
%     only in the simulated plant and not in the baseline.  Theta is drawn
%     once, at the CoM.
%
% \gantryabsorber: 0 draws the baseline only (Section II), 1 adds the grey
% absorber (Section III).
% \gantrydelta: 1 adds an arrow for the absorber displacement relative to the
% payload.
% ===========================================================================
\providecommand{\gantryabsorber}{1}
\providecommand{\gantrydelta}{1}
\begin{tikzpicture}[
    x=0.26mm, y=0.26mm,                      % one column; 0.3mm is 1:1 with
                                             % Garcia Fig. 2
    line width=0.5pt, line cap=round, line join=round,
    >={Triangle[length=3.4pt,width=2.9pt]},
    dim/.style={line width=0.4pt, dash pattern=on 3pt off 1.1pt on 0.4pt off 1.1pt},
    ext/.style={line width=0.35pt},
    cl/.style={gray!55, line width=0.25pt,
               dash pattern=on 4pt off 1pt on 0.3pt off 1pt},
    absc/.style={black!55},                    % absorber colour (not baseline)
    abs/.style={fill=black!15, line width=0.5pt},             % absorber: shaded
    absl/.style={font=\small},                                % k_a, c_a labels
  ]
  % ================ parameters ================
  \def\th{9.25}      % yaw angle Theta [deg]
  \def\Lb{152.5}     % cross arm length
  \def\Wb{24.5}      % cross arm width
  \def\Yp{-18}       % payload position along the cross arm (u); Y = -u
  \def\Lh{26}        % payload width
  \def\Hh{24}        % payload height
  \def\rw{3.7}       % payload wheel radius
  \def\Wa{20}        % absorber width
  \def\Ha{20}        % absorber height
  \def\ga{30}        % spring and damper gap between absorber and payload
  \def\Mw{34.7}      % actuator mass box size
  \pgfmathsetmacro\yE{0.5*\Lb*sin(\th)}       % height of the beam ends
  \pgfmathsetmacro\hb{0.5*\Wb}
  \pgfmathsetmacro\vh{\hb+2*\rw+0.5*\Hh}      % payload centre height (v)
  \pgfmathsetmacro\ua{\Yp-0.5*\Lh-\ga-0.5*\Wa}  % absorber centre (u)
  \pgfmathsetmacro\vY{\vh+0.5*\Hh+30}          % height of the Y dimension line
  \pgfmathsetmacro\xM{0.5*\Lb*cos(\th)+22.3+0.5*\Mw}  % |x| of actuator centres
  % Lane for the d dimension and the F_Y arrow: on the CoM side of the payload
  % and clear of the red Y reference line, so it follows the sign of \Yp.
  \pgfmathsetmacro\ud{ifthenelse(\Yp<0,9,-9)}
  \pgfmathsetmacro\sd{ifthenelse(\Yp<0,1,-1)}

  % ===================== cross arm =====================
  \begin{scope}[rotate=\th]
    \draw (-\Lb/2,-\hb) rectangle (\Lb/2,\hb);
    \draw[cl] (-0.57*\Lb,0) -- (0.58*\Lb,0);
    \coordinate (E1) at ( \Lb/2,0);
    \coordinate (E2) at (-\Lb/2,0);

    % ==== payload m_h on its wheels ====
    \draw[fill=white] (\Yp-\Lh/2,\vh-\Hh/2) rectangle (\Yp+\Lh/2,\vh+\Hh/2);
    \draw (\Yp-7,\hb+\rw) circle (\rw)  (\Yp+7,\hb+\rw) circle (\rw);
    \node at (\Yp,\vh) {\ifnum\gantryabsorber>0 $m_{h,r}$\else $m_h$\fi};

    % ==== payload absorber: mass m_a, spring k_a, damper c_a ====
    % Drawn entirely in grey: not part of the baseline, only of the simulated
    % data-generating plant.
    \ifnum\gantryabsorber>0
    \begin{scope}[absc]
    \pgfmathsetmacro\xa{\ua+0.5*\Wa}          % absorber side attachment
    \pgfmathsetmacro\xh{\Yp-0.5*\Lh}          % payload side attachment
    \draw[abs] (\ua-\Wa/2,\vh-\Ha/2) rectangle (\ua+\Wa/2,\vh+\Ha/2);
    \node at (\ua,\vh) {$m_a$};
    \pgfmathsetmacro\xc{0.5*(\xa+\xh)}        % centre of the gap
    % spring k_a on top
    \draw[decorate, decoration={zigzag, segment length=4, amplitude=3.5,
          pre length=5, post length=5}] (\xa,\vh+6) -- (\xh,\vh+6);
    \node[absl, anchor=south, inner sep=1.5pt] at (\xc,\vh+10) {$k_a$};
    % damper c_a below: rod, filled piston, open cylinder, rod
    \draw (\xa,\vh-6) -- (\xc+1,\vh-6);
    \fill (\xc-1.4,\vh-8.6) rectangle (\xc+0.8,\vh-3.4);
    \draw (\xc-4,\vh-2) -- (\xc+4,\vh-2) -- (\xc+4,\vh-10) -- (\xc-4,\vh-10);
    \draw (\xc+4,\vh-6) -- (\xh,\vh-6);
    \node[absl, anchor=north, inner sep=0.5pt]        % tucked up under the
          at (\xc,\vh-10.9) {$c_a$};                   % cylinder, clear of the beam
    \ifnum\gantrydelta>0                      % displacement relative to m_h
      \draw[->] (\ua,\vh+\Ha/2+4.5) -- ++(-14,0);
      \node[anchor=south east, inner sep=1pt]
            at (\ua-3,\vh+\Ha/2+5.5) {$q_a$};
    \fi
    \end{scope}
    \fi

    % ==== payload force F_Y ====
    % With the absorber on, the label carries the rail friction (grey, not in
    % the baseline) and is anchored at the arrow tail so it grows away from it
    % (Jasper v1 p. 9, "indicate the friction too"; Dirk 2026-10-24).
    \draw[->] ({\ud+\sd*11},\vh+4) -- ({\Yp+\sd*\Lh/2},\vh+4);
    \ifnum\gantryabsorber>0
      \ifdim\Yp pt<0pt \def\FYanchor{west}\else \def\FYanchor{east}\fi
      \node[inner sep=2pt, anchor=\FYanchor]
            at ({\ud+\sd*12},\vh+4)
            {$F_Y$\textcolor{black!55}{$\,-\,F_{\mathrm{fr},Y}$}};
    \else
      \node[inner sep=2pt] at ({\ud+\sd*20},\vh+4) {$F_Y$};
    \fi

    % ==== d : centreline to payload centre of mass ====
    \draw[<->] (\ud,0) -- (\ud,\vh);
    \node[inner sep=2pt] at ({\ud+\sd*4.5},0.78*\vh) {$d$};  % away from CoM

    % ==== Y : CoM to payload (red) ====
    \begin{scope}[red]
      \draw[dim] (0,\hb+2) -- (0,\vY+6)               % extension lines at the
                 (\Yp,\vh+\Hh/2+2) -- (\Yp,\vY+6);    % CoM and at the payload
      \draw[->] (0,\vY) -- (\Yp,\vY);
      \node[anchor=south] at ({0.5*\Yp},\vY+1) {$Y$};
    \end{scope}

    % ==== flexible joints and L_b ====
    \draw[dim] ( \Lb/2,-\hb-2) -- ++(0,-24)   (-\Lb/2,-\hb-2) -- ++(0,-24);
    \draw[<->] (-\Lb/2,-\hb-18) -- node[pos=0.72, below=1pt] {$L_b$}
               ( \Lb/2,-\hb-18);
    \node[anchor=north west, inner sep=2.5pt] at ( \Lb/2,-\hb-19) {$k_{b1},\,c_{b1}$};
    \node[anchor=north east, inner sep=2.5pt] at (-\Lb/2,-\hb-19) {$k_{b2},\,c_{b2}$};
  \end{scope}

  % ===================== generalized coordinates X, Theta (blue) =============
  \begin{scope}[blue]
    \draw[->] (0,-5.2) -- (0,-26);
    \node[anchor=west, inner sep=2pt] at (1.5,-21) {$X$};
    \draw (4.6,0) -- (36,0);
    \draw (32,0) arc[start angle=0, end angle=\th, radius=32];
    \node[anchor=west, inner sep=2.5pt] at (37,{36*sin(0.5*\th)}) {$\Theta$};
  \end{scope}
  % centre of mass m_b
  \fill[white] (0,0) circle (4.3);
  \fill (0,0) -- (4.3,0) arc[start angle=0, end angle=-90, radius=4.3] -- cycle
        (0,0) -- (-4.3,0) arc[start angle=180, end angle=90, radius=4.3] -- cycle;
  \draw (0,0) circle (4.3);
  \node[anchor=north east, inner sep=1pt] at (-3.4,-3.4) {$m_b$};

  % ===================== actuators, guides, forces =====================
  % \s = +1 : actuator 1 (right) ; \s = -1 : actuator 2 (left)
  \foreach \s/\i in {1/1,-1/2}{
    \pgfmathsetmacro\y{\s*\yE}                 % row height (symmetric)
    \coordinate (M\i) at (\s*\xM,\y);
    \draw (\s*\xM-\Mw/2,\y-\Mw/2) rectangle (\s*\xM+\Mw/2,\y+\Mw/2);
    \node at (M\i) {$m_\i$};

    % flexible joint: spiral spring between beam end and actuator
    \pgfmathsetmacro\xs{\s*(0.5*\Lb*cos(\th)+11)}
    \draw (\xs,\y) -- (E\i);
    \draw[line width=0.4pt] plot[domain=0:540, samples=100, smooth]  % 1.5 turns
          ({\xs+(3.4+3.9*\x/540)*cos(90*(1-\s)-\s*(\x-540))},
           {\y +(3.4+3.9*\x/540)*sin(90*(1-\s)-\s*(\x-540))});
    \draw (\xs+\s*7.3,\y) -- (\s*\xM-\s*\Mw/2,\y);

    % linear guide: two rollers against a hatched wall
    \pgfmathsetmacro\xb{\s*(\xM+\Mw/2+5.35)}
    \draw (\xb,\y+11.6) circle (5.35)  (\xb,\y-11.6) circle (5.35);
    \draw (\xb+\s*5.6,\y-20) -- (\xb+\s*5.6,\y+20);
    \foreach \k in {0,...,5}
      \draw[ext] (\xb+\s*5.6,\y+18-7.5*\k) -- ++(\s*6.5,-6.5);

    % actuator force
    \draw[->] (\s*\xM,\y+\Mw/2+30) -- (\s*\xM,\y+\Mw/2);
    \ifnum\gantryabsorber>0                   % rail friction, grey; the label
      \ifnum\i=2 \def\FXanchor{south east}\def\FXdx{\Mw/2}% grows outward,
      \else \def\FXanchor{south}\def\FXdx{0}\fi          % clear of q_a
      \node[anchor=\FXanchor] at (\s*\xM+\FXdx,\y+\Mw/2+31)
            {$F_{X\i}$\textcolor{black!55}{$\,-\,F_{\mathrm{fr},X_\i}$}};
    \else
      \node[anchor=south] at (\s*\xM,\y+\Mw/2+31) {$F_{X\i}$};
    \fi

    % measured position X_i (red), drawn from the carriage so that it does
    % not start at the fixed world (feedback Js24, 2026-10-18); was a tick
    % and arrow outside the wall hatching
    \pgfmathsetmacro\xr{\s*(\xM+\Mw/2-1.5)}   % outer edge of the carriage,
                                              % clear of the k_b, c_b labels
    \draw[red,->] (\xr,\y-\Mw/2) -- (\xr,\y-\Mw/2-22);
    \node[red, anchor=north] at (\xr,\y-\Mw/2-23) {$X_\i$};
  }

  % ===================== inertial frame (top view) =====================
  \begin{scope}[shift={(72,-58)}]
    \draw[->] (0,-3.2) -- (0,-19);          % x points down
    \draw[->] (-3.2,0) -- (-19,0);          % y points left
    \draw[fill=white] (0,0) circle (3.2);   % z into the page: circle with a cross
    \draw[ext] (-2.26,-2.26) -- (2.26,2.26)  (-2.26,2.26) -- (2.26,-2.26);
    \node[anchor=east, inner sep=2pt]  at (-19,0.3) {$y$};
    \node[anchor=west, inner sep=2pt]  at (3.6,0.8) {$z$};
    \node[anchor=north, inner sep=2pt] at (0,-20)  {$x$};
  \end{scope}
\end{tikzpicture}
}
  \caption{Simulated gantry with a payload-mounted absorber. The light-gray
  elements $(m_a,k_a,c_a,q_a)$ are absent from the baseline.}
  \label{fig:gantry_absorber}
\end{figure}

The friction acts on each rail against its actuator force.
With it, \eqref{eq:eom} becomes
\begin{subequations}\label{eq:setup_truth}
\begin{equation}
\begin{aligned}
 M(Y)\ddot q+C\dot q+Kq&=P\bigl(\uact-F_{\mathrm{fr}}(P^\top\dot q)\bigr),\\
 F_{\mathrm{fr},i}(v_i)&=c_{c,i}\tanh\bigl(v_i/v_{\mathrm{fr}}\bigr),
\end{aligned}
 \label{eq:setup_fric}
\end{equation}
where $P^\top\dot q$ are the rail velocities, $i\in\{X_1,X_2,Y\}$, and
$v_{\mathrm{fr}}$ is the velocity scale of the smoothing.
For $\lvert v_i\rvert\gg v_{\mathrm{fr}}$, $F_{\mathrm{fr},i}$ is the
Coulomb law $c_{c,i}\operatorname{sign}(v_i)$ of~\cite{garcia2013model}.
With $v_{\mathrm{fr}}=\SetupFricVzero$~mm/s, this holds on most samples of
the records with multisine (Table~\ref{tab:setup_checks}).
At standstill the law acts as a stiff viscous damper, so the axes settle
without a limit cycle.

The absorber adds the displacement $q_a$ of its mass relative to the head
along $Y$.
With $q_{\mathrm e}=\begin{bmatrix}q^\top&q_a\end{bmatrix}^\top$, the
system is
\begin{equation}
 M_{\mathrm e}(Y,q_a)\,\ddot q_{\mathrm e}+C_{\mathrm e}\dot q_{\mathrm e}+K_{\mathrm e}q_{\mathrm e}
 =\begin{bmatrix}P\bigl(\uact-F_{\mathrm{fr}}(P^\top\dot q)\bigr)\\0\end{bmatrix},
 \label{eq:setup_abs}
\end{equation}
\end{subequations}
with
{\footnotesize
\setlength{\arraycolsep}{3pt}
\begin{equation*}
\begin{gathered}
 M_{\mathrm e}(Y,q_a)=\begin{bmatrix}M(Y)+m_aE(Y,\ell)&m_ab_a\\m_ab_a^\top&m_a\end{bmatrix},\\
 E(Y,\ell)=\begin{bmatrix}0&-\ell&0\\-\ell&\ell(2Y+\ell)&0\\0&0&0\end{bmatrix},\\
 b_a=\begin{bmatrix}0&-d&1\end{bmatrix}^\top,\quad
 C_{\mathrm e}=\operatorname{diag}(C,c_a),\quad
 K_{\mathrm e}=\operatorname{diag}(K,k_a),
\end{gathered}
\end{equation*}}
where $M(Y)$ is \eqref{eq:mass_matrix} with $m_h$ the total head mass, of
which $m_a$ moves relative to the head, $\ell=L_0+q_a$, and $L_0$ is the
offset of the absorber along $Y$.
The absorber enters through the inertia, by $E$ and $b_a$, and through its
own spring and damper.
The system is integrated by fixed-step RK4 at $f_{\mathrm s}=20$~kHz.

Table~\ref{tab:setup_truth} gives the parameters that the baseline lacks.
The baseline parameters are those of Appendix~\ref{app:params}.
Seen from the head, $f_a$ is an anti-resonance.
It and the coupled mode lie above the crossover of the loop
(Section~\ref{sec:setup_data}), so the feedback does not suppress them.
\begin{table}[t]
\caption{Parameters that the baseline lacks.}
\label{tab:setup_truth}
\centering
\footnotesize
\setlength{\tabcolsep}{4pt}
\begin{tabular}{@{}lll@{}}
\toprule
Parameter & Value & Basis\\
\midrule
$c_c$ ($X_1$, $X_2$, $Y$) & 16.80, 18.35, 11.60~N & \cite[Table~I]{garcia2013model}\\
$v_{\mathrm{fr}}$ & \SetupFricVzero~mm/s & chosen\\
$m_a$ & $0.5\,m_h$ & chosen\\
$L_0$ & 0.10~m & chosen\\
$f_a$ & 150~Hz & chosen\\
$\zeta_a$ & 0.03 & jointed metal\\
$k_a$, $c_a$ & $m_a(2\pi f_a)^2$, $2\zeta_a\sqrt{k_am_a}$ & derived\\
Coupled mode & $f_a\sqrt{1+m_a/(m_h-m_a)}\approx212$~Hz & derived\\
\bottomrule
\end{tabular}
\end{table}
\todo{Missing: reason for $m_a=0.5\,m_h$, $L_0=0.10$~m and $f_a=150$~Hz (inherited from \texttt{tdg\_config.m}, no reason on record), and a source for $\zeta_a$. Candidate: $m_a$ large enough that the mode carries a learnable share of the baseline error (D-170, D-188); $\zeta_a$ the low end of the damping range of jointed metal structures (D-188, web source \texttt{jpe\_structural\_damping}, no cite key).}

The records are closed loop.
Each stage axis $j$ has one controller $K_{1,j}$ with 100~Hz bandwidth,
designed on the rigid baseline at $Y=0$ (Appendix~\ref{app:setup_system}).
$K_1$ is the same in every record and is the known controller of
Section~\ref{sec:aug_loop}.
The feedforward $u^{\mathrm{ff}}$ of \eqref{eq:aug_two_loops} is the
injected excitation force only.

The position measurement carries noise, $y=\qact+v$.
The components of $v$ are independent and coloured.
Their standstill servo error matches the level and spectral shape of the
industrial machine, without its resonances.
Every record also exists without noise, denoted $y^{\circ}$.
The signal-to-noise ratio of the servo error is \SetupSnrMinX{} to
\SetupSnrMaxX~dB on the $X$ rails and \SetupSnrMinY{} to \SetupSnrMaxY~dB
on $Y$.
\todo{Missing: what may be named about the machine and its noise. Candidate: ask Jasper or Dragan; if allowed, list $f_{v,j}$, $\lambda_j$, $\sigma_{v,j}$ in Appendix~\ref{app:setup_system}.}

\subsection{Records}\label{sec:setup_data}
% Sources: D-219, D-226 (band), D-221 (split), D-258 (twins), README validation story;
% Code/Setup band_edge.py, band_frf.py, coverage.py, rate.py, filter_leak.py, loop_memory.py.
\ifdraft
\begin{itemize}
\item A record: 12~s in closed loop with $K_1$, one of five motion types, a multisine added except on the operational moves; noisy and noise-free
\item The multisine band starts at the highest crossover of the loop and holds the absorber features of the simulated system (simulation only)
\item Split 18 / 6 / 11; validation only selects; every reported number on the test records
\item Test records by layer: interpolation, operational motion, extrapolation in one setting, stress record E6, multisine twins for evaluation only
\item Model rate, anti-alias filter and window checked on the simulated loop (appendix)
\end{itemize}
\fi
Each record is a closed-loop simulation of 12~s at $f_{\mathrm s}$ with
$K_1$.
It follows one of five motion types, from standstill to sine-shaped moves
of machine operation (Appendix~\ref{app:setup_records}).
Except on the sine-shaped moves, a multisine force is injected in each
channel $X$, $\Theta$, $Y$ of the generalised force.
Every record exists with and without noise.

The motion profiles excite the low frequencies.
The multisine has to excite the absorber, but below the crossover of the
loop the feedback suppresses an injected force.
Its band of 106 to 297~Hz therefore starts at the highest crossover of the
loop with $K_1$ over the axes and the training positions,
\SetupCrossMax~Hz.
Its upper edge keeps 90\,\% of the difference between the frequency
responses of the system and the baseline above the crossover
(Appendix~\ref{app:setup_records}).
On the simulated system, the band holds the anti-resonance, the coupled
mode and the closed-loop peak at \SetupClPeak~Hz at all five training
positions (Fig.~\ref{fig:setup_band}).
This placement uses the known system and is available only in simulation.
On a machine, a measured frequency response plays this role.

Of the records, 18 are used for training, 6 for validation and 11 for
testing (Fig.~\ref{fig:setup_coverage}).
Training covers $\lvert Y\rvert\le\SetupTrainY$~m, accelerations up to
\SetupTrainAcc\,\% and velocities up to \SetupTrainVel\,\% of the machine
limits.
The validation records use new positions and new random draws.
They only select the checkpoints and the size of the black box.
Every reported number is computed on the test records of
Table~\ref{tab:setup_records}, which test interpolation, operational motion
and range extrapolation separately.
Each extrapolation record changes one design setting against an
interpolation record.
E6 differs from I3 in more than one property and is reported as a stress
record.
\todo{Missing: whether E6 (a measured machine move) and the relative limits may be described this way. Candidate: ask Jasper or Dragan.}
\begin{table}[t]
\caption{Test records. The twins add the training multisine, at 0.30 of
its amplitude for the pair I6, E4.}
\label{tab:setup_records}
\centering
\footnotesize
\setlength{\tabcolsep}{3pt}
\begin{tabular}{@{}llll@{}}
\toprule
Record & Motion & Multisine & Tests\\
\midrule
I1 & training type & yes & new phases and setpoints\\
I2 & standstill & yes & new position\\
I3 to I6 & machine-like & no & operational motion\\
E1 & as I2 & yes & position, $\lvert Y\rvert=\SetupEY$~m\\
E2 & as I3 & no & stroke of $Y$\\
E3 & as I5 & no & acceleration limit, \SetupEAcc\,\%\\
E4 & as I6 & no & velocity limit, \SetupEVel\,\%\\
E6 & machine move & no & stress, \SetupESixAcc\,\% acceleration\\
I3M to E6M & twins & yes & absorber band on every pair\\
\bottomrule
\end{tabular}
\end{table}

The models run at $T_s=0.25$~ms (4~kHz) on training windows of 0.1~s.
The reduction of the records to this rate, its anti-alias filter and the
window length are checked on the simulated loop in
Table~\ref{tab:setup_checks}.

\subsection{Compared models}\label{sec:setup_models}
% Sources: THESIS-RESULTS.md "Presentation agreed 2026-10-16"; D-191, D-243 (detuned start);
% D-242, thesis_r8_grid.conf (black box); D-259 (single seed).
\ifdraft
\begin{itemize}
\item Each comparison answers one aspect of the research question
\item Grey-box models: same records, settings, budget and checkpoint rule; PS start of the added states; 10\,\% detuned start of the combinations; one seed
\item The runs that justify $n_{\bar x}$, PS and the start are in the appendix (?: label)
\item Black box: SUBNET of the closed loop, no physics, no controller, sized on the validation records
\end{itemize}
\fi
Each comparison between the models of Table~\ref{tab:setup_models}
answers one aspect of the research question (last column).
The pairs are $\mathcal M_{\mathrm{LPV}}$ and $\mathcal M_{\mathrm{LTI}}$
for the baseline, the final model $\mathcal M_{\mathrm{OBC}}$ with
$\mathcal M_{\mathrm{LPV}}$ and $\mathcal M_{\mathrm{BB}}$ for the
augmentation, and $\mathcal M_{\mathrm{OBC}}$ with $\mathcal M_{\mathrm U}$
for the orthogonal construction.
The extrapolation records compare the final model with the black box
(Aspect~4).
$\mathcal M_{\mathrm{LPV}}$ and $\mathcal M_{\mathrm{LTI}}$ are the baseline
at the true parameters without training.
$\mathcal M_{\mathrm{LTI}}$ freezes $M(Y)$ at $Y=0$, the middle of the
stroke.
\begin{table}[t]
\caption{Compared models. $\theta^0_{\mathrm{base}}$: start of the ten
combinations; final model: $\mathcal M_{\mathrm{OBC}}$.}
\label{tab:setup_models}
\centering
\footnotesize
\begin{tabular}{@{}lllll@{}}
\toprule
Model & $n_{\bar x}$ & OBC & $\theta^0_{\mathrm{base}}$ & Aspect\\
\midrule
$\mathcal M_{\mathrm{LPV}}$ & 0 & no & true, untrained & 1, 2\\
$\mathcal M_{\mathrm{LTI}}$ & 0 & no & true at $Y=0$, untrained & 1\\
$\mathcal M_{\mathrm{U}}$ & 2 & no & detuned & 3\\
$\mathcal M_{\mathrm{OBC}}$ & 2 & yes & detuned & 2, 3, 4\\
$\mathcal M_{\mathrm{BB}}$ & \multicolumn{3}{l}{black box, Appendix~\ref{app:setup_bb}} & 2, 4\\
\bottomrule
\end{tabular}
\end{table}
\todo{Missing: model notation. Candidate: the calligraphic names of the table (README open conflict on model names; same \texttt{\textbackslash todo} as in Section~\ref{sec:obc}).}

$\mathcal M_{\mathrm U}$ and $\mathcal M_{\mathrm{OBC}}$ are identified by
the criterion of Section~\ref{sec:aug_closed_loop} on the same training
records, with the same settings, budget and checkpoint rule
(Table~\ref{tab:hyper}).
Both start the added states by PS (Section~\ref{sec:aug_ps2}).
Both start the combinations at
$\theta^0_{\mathrm{base},j}=(1+0.1\,s_j)\,\theta^\star_{\mathrm{base},j}$,
with signs $s_j\in\{-1,1\}$ that are the same in every run.
The final model therefore holds no true parameter value, like the black
box.
The runs that justify $n_{\bar x}$, PS and the start are in the appendix.
Every model is trained once, with one seed (compute limit), so the spread
over seeds is not measured and small differences between models are not
resolved.
\todo{Missing: the label of the appendix with the choice-justifying runs, and whether they are compared on the validation records. Candidate: a new appendix section on the validation records (THESIS-RESULTS, recommendation of 2026-10-16, not yet confirmed).}

The black box $\mathcal M_{\mathrm{BB}}$ is a SUBNET
model~\cite{beintema2023subnet}, a state-space model with neural maps
identified on simulation windows.
It models the closed loop from the reference and the injected force to the
servo error (Appendix~\ref{app:setup_bb}).
It uses no physical knowledge and no controller.
Its state count $n_x$, network widths and learning rate are selected on the
validation records from 11 configurations, with the data and windows of
the grey-box models.
\revnote{In the current search, the selected configuration ($n_x=16$, width
8) had its best validation error at the limit of 20000 updates; the final
retraining to a validation plateau has not run yet.}

\subsection{Metrics}\label{sec:setup_metrics}
% Sources: D-256 (NRMS_e), THESIS-RESULTS Sec. 2; D-259 (m_Delta scale, revnotes); Section IV
% eq:obc_eps, eq:obc_recovery, eq:obc_overlap; Code/Evaluation eval_overlap.py.
\ifdraft
\begin{itemize}
\item Every model is simulated over whole test records: grey-box models in closed loop with $K_1$, the black box in free run
\item Primary metric: the simulation error relative to the servo error, $\mathrm{NRMS}_e$, per stage axis; beside it the RMS in metres and the noise floor
\item Parameter error per combination against the true values; $m_\Delta$ scaled by the mean actuator mass; simulation only
\item The two predictions of Section~\ref{sec:obc}: the network stays out of the parameter directions, and the remaining parameter error is the first-order bias
\end{itemize}
\fi
Each grey-box model is simulated in closed loop with $K_1$
(Section~\ref{sec:aug_loop}) over each whole test record, from the encoder
estimate at the record start.
The black box is simulated in free run over the same records.

The reference is given, so a model has to predict the servo error $y-r$.
The primary metric is therefore the RMS of the simulation error per stage
axis $j$ relative to the standard deviation of the servo error.
Beside it, the RMS in metres and the noise floor of the record are
reported:
\begin{equation}
\begin{aligned}
 \mathrm{RMS}_j&=\Bigl(\frac{1}{N-n}\sum_{k=n}^{N-1}(\hat y_{j,k}-y_{j,k})^2\Bigr)^{1/2},\\
 \mathrm{NRMS}_{e,j}&=\frac{\mathrm{RMS}_j}{\sigma_{e,j}},\\
 \mathrm{RMS}^{\mathrm{fl}}_j&=\Bigl(\frac{1}{N-n}\sum_{k=n}^{N-1}(y_{j,k}-y^{\circ}_{j,k})^2\Bigr)^{1/2},
\end{aligned}
\label{eq:setup_rms}
\end{equation}
where $N$ is the record length, $n$ the encoder window, and
$\sigma_{e,j}$ is the standard deviation of the servo error $y_j-r_j$ over
the same samples.
Since $\hat y-y=(\hat y-r)-(y-r)$, $\mathrm{NRMS}_{e,j}$ is the NRMS
of~\cite[Eq.~(43)]{beintema2023subnet} with the servo error as output.
A value of 100\,\% means no better than $\hat y=r$.
The floor is the error of an exact noise-free prediction.
\revnote{The current evaluation reports the RMS in metres per record and
axis; $\mathrm{NRMS}_{e,j}$ is computed from these outside the evaluation
code.}
\revnote{Start index of \eqref{eq:setup_rms}: the window $\chi_\tau$ holds
the $n+1$ samples up to $\tau-1$ (\eqref{eq:aug_encoder}), so the first
window covers samples $0$ to $n$ and the first prediction is at $k=n+1$.
Following~\cite[Eq.~(43)]{beintema2023subnet}, which leaves out the outputs
that fed the encoder, the sums run from $k=n+1$ with $1/(N-n-1)$; the code
of the final runs scores from that sample. $\sigma_{e,j}$ is the population
standard deviation over the same samples (D-262).}

The combinations are compared with their true values
$\theta^\star_{\mathrm{base}}$ by
$\delta_j=\lvert\theta_{\mathrm{base},j}-\theta^\star_{\mathrm{base},j}\rvert/\lvert\theta^\star_{\mathrm{base},j}\rvert$.
For $m_\Delta=m_1-m_2$, the denominator is the mean actuator mass
$(m_1+m_2)/2$ instead.
Its true value of $-0.5$~kg is close to zero, while it acts in $M(Y)$
beside the actuator masses.
This uses the true parameters and is available only in simulation.

Section~\ref{sec:obc} predicts two properties of
$\mathcal M_{\mathrm{OBC}}$.
Both are measured on the one-step tuples of I1 to I6 at $\bar x=0$.
The first is that the network stays out of the parameter directions.
It is measured by the overlap $\rho$ of \eqref{eq:obc_overlap}, with $\Phi$
at the trained coordinates, for $F^{\mathrm n}_{\mathrm{aug}}$ of
$\mathcal M_{\mathrm U}$ and for $\tilde F^{\mathrm n}_{\mathrm{aug}}$ of
\eqref{eq:obc_corr} for $\mathcal M_{\mathrm{OBC}}$.
The second is that the remaining parameter error is the first-order bias of
\eqref{eq:obc_recovery}.
Its first line, with $\Xi$ of \eqref{eq:obc_eps} and $\Phi$ at
$\beta^\star$, predicts the $\delta_j$ of the trained
$\mathcal M_{\mathrm{OBC}}$.
The truth overlap $\rho(\Xi)$ measures how far the missing dynamics lie in
the parameter directions.
\revnote{The predicted first-order bias is not yet computed for the current
runs.}
\revnote{Final runs (D-268, Dirk 2026-10-16): the overlaps and the predicted
first-order bias are computed on three tuple stacks, training, validation and
test (I1 to I6), so the sentence ``Both are measured on the one-step tuples of
I1 to I6'' changes. The training stack, where $\theta_{\mathrm{aux}}$ is
formed, gives the prediction of \eqref{eq:obc_recovery}; validation and test
show how it carries over. On the training stack the network overlap of
$\mathcal M_{\mathrm{OBC}}$ is zero by construction, so that overlap is
informative on validation and test. The prediction holds for the one-step
least-squares fit at $\bar x=0$, while training minimises the closed-loop loss
with active added states (Section~\ref{sec:obc}); the agreement criterion
between predicted and observed $\delta_j$ is to be stated before the results.}
